\PassOptionsToPackage{sols}{handout}
\documentclass[main]{subfiles}
\setcounterrefbeforedot{chapter}{m-ws7-1}
\setcounterrefafterdot{section}{m-ws7-1}
\addtocounter{section}{-1}
\usepackage{solutions}

\begin{document}

\ifSubfilesClassLoaded{%
  \chaptermark{\chapterseven}
}{}

\section{Introduction to Differential Equations} \label{ws7-1}

\begin{problemonly}
    \begin{framed}

\textbf{Key Ideas}:

\begin{itemize}[topsep=0pt,partopsep=0pt,parsep=8pt,itemsep=0pt]

\item You should understand what we mean by solutions of differential
  equations and the \uline{general solution} of a differential equation.

\item You should be able to determine whether a given function is a
  solution of a given differential equation.

\item You should be able to find the general solution of a differential
  equation of the form $\dfrac{dy}{dt} = f(t)$.

\item You should be comfortable with the differential equation
  $\frac{dy}{dt} = ky$ where $k$ is a constant.  You should understand
  that this differential equation says that $y$ changes at a rate
  proportional to itself and know the general solution of this
  differential equation.

\end{itemize}
\end{framed}
\end{problemonly}

A derivative gives
\begin{itemize}[topsep=0pt]
    \item an instantaneous rate of change
    \item slope at a point
\end{itemize}
The former interpretation allows us to model situations by describing rates of change. It is often simpler to describe a phenomenon from the perspective of rates of change. The latter helps us think about solutions.

  \begin{framed}
    A \uline{differential equation} is an equation that involves an unknown function and one or more of its derivatives.
  \end{framed}



\begin{enumerate}
\item  \label{diffeq-modeling-bacteria-1}

  The population of a bacteria colony, under ideal conditions, increases at a rate proportional to
  the population size. For this particular model, we will assume ideal growing conditions. (More sophisticated models should include competition for limited resources.)
  \begin{enumerate}
  \item
    \begin{problem}[0.5in]
      If $P(t)$ is the population size at time $t$, write a differential
      equation that expresses the situation.      
    \end{problem}

  \begin{framed}
     {$A$ is proportional to $B$} translates mathematically to $A = kB$ for $k$ a constant.
  \end{framed}
    \begin{solution}
      We can translate the information we're given into an equation: 
      \[ \underbrace{\text{The rate at which the population
          increases}}_{\frac{dP}{dt}} \text{ } \underbrace{\text{is
          proportional to}}_{{} = k \text{ times}} \text{ }
      \underbrace{\text{the population size}}_{P}.\] That is,
      $\boxed{\frac{dP}{dt} = k P}$. 
    \end{solution}

  \item
    \begin{problem}[0.4in]
      When the colony has a population of 500 bacteria, the colony is
      growing at a rate of 100 bacteria per hour.  Incorporate this into
      your differential equation.
    \end{problem}

    \begin{solution}
      We are told that, when $P = 500$, $\frac{dP}{dt} = 100$.  Plugging
      this into our differential equation $\frac{dP}{dt} = kP$, $100 = 500
      k$, so $k = 0.2$.  Therefore, our differential equation becomes
      $\boxed{\frac{dP}{dt} = 0.2 P}$.
    \end{solution}
\item
    \begin{problem}[\vfill]
      Does it take more than an hour, less than an hour, or exactly one
      hour for the population to grow from 500 to 600 bacteria?  Explain
      how you can be sure.
    \end{problem}

    \begin{solution}
      We were told that, when the colony has a population of 500 bacteria,
      the colony is growing at an instantaneous rate of 100 bacteria per
      hour.  As the colony grows, its rate of growth also grows (because
      its rate of growth is proportional to the size of the colony).
      Thus, the colony will immediately start growing faster than 100
      bacteria per hour, so it will take \fbox{less than one hour} for the
      population to grow to 600.
    \end{solution}
 
  \item
    \begin{problem}[\newpage]
      If the colony has an initial population of 200 bacteria, we should be able to find a
      formula for $P(t)$. We will set this aside for a moment and look at a similar, but simpler example to understand what it means to be a solution to a differential equation and then will return to this problem.
    \end{problem}

    \begin{solution}
      The general solution to the differential equation $\frac{dP}{dt} =
      0.2 P$ is $P(t) = C e^{0.2 t}$.  Knowing the initial population is
      200 says that $P(0) = 200$, so $C e^{0.2 (0)} = 200$, which tells us
      that $C = 200$.  So, $P(t) = \boxed{200 e^{0.2 t}}$.
    \end{solution}

  \end{enumerate}
  
\item 
\note{This problem can be used to show students they can check to see if
a function is or is not a solution to a differential equation.}

  Let's look at the differential equation $\dfrac{dy}{dt} = y$.

  \begin{enumerate}
  \item

    \begin{problem}[0.4in]
      This differential equation is telling us to look for
      \cfitb{0.5in}{$y$} as a function of \cfitb{0.5in}{$t$}.  Write the
      differential equation in words.
    \end{problem}

    \begin{solution}
      The fact that we're looking at a derivative $\frac{dy}{dt}$ tells us
      that we're thinking of $y$ as a function of $t$.

      In words, the equation $\dfrac{dy}{dt} = y$ says, ``the
      derivative of the function $y(t)$ is equal to itself''.
    \end{solution}

  \item

    \begin{problem}
      Which of the following functions are solutions of the differential
      equation $\dfrac{dy}{dt} = y$?
      
        \begin{enumerate}[label=\protect{\maybecircle{\Alph*.}}, itemsep=0.2cm]
        \plainitem $y(t)=\dfrac{t^2}{2}+7$
        \plainitem $y(t)=\dfrac{(t+1)^2}{2}$
        \circleitem $y(t)=e^t$
        \plainitem $y(t)=e^t+5$
        \circleitem $y(t)=7e^t$
        \end{enumerate}      
            
    \end{problem}
  \begin{framed}
    A function that satisfies the differential equation is a particular solution to the differential equation. Often a \textbf{ particular solution}  is specified by an \textbf{initial condition}. \\

  \end{framed}
    \begin{solution}
      For each function $y(t)$, we'd like to determine whether the
      derivative of the function $y(t)$ is equal to itself (that's exactly what we said in
      {{\#1}}).  Let's look at each one:
      \begin{enumerate}[label=\Alph*.]
      \item If $y(t)=\dfrac{t^2}{2}+7$, the derivative $\frac{dy}{dt}$ is
        equal to $t$, which is not equal to $y(t)$.
      \item If $y(t)=\dfrac{(t+1)^2}{2}$, then $\frac{dy}{dt} = t+1$,
        which is not equal to $y(t)$.
      \item If $y(t)=e^t$, then
        $\frac{dy}{dt} = e^t$, which is equal to $y(t)$.
      \item If $y(t)=e^t+5$, then
        $\frac{dy}{dt} = -e^t$, which is not equal to $y(t)$.
      \item If $y(t)=7e^t$, then $\frac{dy}{dt} = 7e^t$, which is equal
      to $y(t)$.
      \end{enumerate}
    \end{solution}

  \item

    \begin{problem}[\vfill]
      Can you guess any more solutions of the differential equation? 
    \end{problem}

    \begin{solution}
      Two of the solutions we saw were $y(t)=e^t$ and $y(t)=7e^t$, so we might guess that any function of the form
      $y(t) = Ce^t$ where $C$ is a constant is a solution of the
      differential equation. 

      It's easy to check whether our guess is correct, so let's do that.
      If $y(t) = Ce^t$, the derivative $\frac{dy}{dt}$ is
      $Ce^t$, which is exactly what we want.
    \end{solution}

  \end{enumerate}
\item 
\begin{problem}[\vfill]
We want to find solutions to the differential equation $\dfrac{dy}{dt} = ky$ where $k$ is a constant.
\begin{enumerate}

  \item
    \note{We talked last semester about the fact that exponential
      functions are functions for which the rate of growth is proportional
      to the function itself. (If that point was not made strongly, make that point now!)  I won't attempt to justify rigorously that
      $C e^{2t}$ are the only solutions here.}

    \begin{problem}[\vfill]
      Can you think of any solutions of the differential equation
      $\dfrac{dP}{dt} = 2P$?  Can you find the general solution?
    \end{problem}

    \begin{solution}
      We're looking for functions $P(t)$ whose derivative is twice the
      function.  These are exactly functions of the form $\boxed{P(t) = C
        e^{2t}}$.
    \end{solution}

    \item 
    \begin{problem}[\vfill]
        What are some solutions to the differential equation 
        $\dfrac{dP}{dt}=kP$, where $k$ is a constant? Can you find the
        general solution?
    \end{problem}

     \begin{solution}
      We're looking for functions $P(t)$ whose derivative is $k$ times the
      function.  These are exactly functions of the form $\boxed{P(t) = C
        e^{kt}}$.
    \end{solution}

      \end{enumerate}
\end{problem}



\end{enumerate}
\begin{framed}
A  general solution to a differential equation involves an arbitrary constant $C$ (just one constant if the differential equation has first derivatives but no higher order derivatives). The \textbf{general solution} to a differential equation  is a family of functions that are solutions to the differential equation. Any particular solution can be expressed in this form for some constant  $C$.
\end{framed}

\clearpage
  \begin{framed}
    \textbf{Very important fact.}  If $k$ is a constant, then the general
    solution of $\frac{dy}{dt}=ky$ is $y(t) = C
    e^{kt}$.
  \end{framed}


%\note{The next two problems are about modeling with this differential
 % equation.}

\begin{enumerate}[resume]
    \item
    \begin{enumerate}
      
  \item
    \begin{problem}[\stretch{0.8}]
    
      Return to problem \#1(d). If the colony has an initial population of 200 bacteria, find a
      formula for $P(t)$. 
    \end{problem}

    \begin{solution}
      The general solution to the differential equation $\frac{dP}{dt} =
      0.2 P$ is $P(t) = C e^{0.2 t}$.  Knowing the initial population is
      200 says that $P(0) = 200$, so $C e^{0.2 (0)} = 200$, which tells us
      that $C = 200$.  So, $P(t) = \boxed{200 e^{0.2 t}}$.
    \end{solution}
    
  \item
    \begin{problem}[\stretch{1}]
      How long does it take the population to double in size?  To
      quadruple? 
    \end{problem}

    \begin{solution}
      Since the population starts out with 200 bacteria, we're looking for
      the time at which the population reaches 400 bacteria.  Since the
      population at time $t$ is $200 e^{0.2 t}$, we solve
      \begin{align*}
        200 e^{0.2 t} &= 400 \\
        e^{0.2 t} &= 2 \\
        0.2 t &= \ln 2 \\
        t &= 5 \ln 2
      \end{align*}
      So, it takes $\boxed{5 \ln 2}$ hours for the population to double
      (this is about $3.47$ hours).

      Quadrupling is the same as doubling twice, so it takes $2 \cdot 5
      \ln 2 = \boxed{10 \ln 2}$ hours for the population to quadruple. 
    \end{solution}
    
    \end{enumerate}

\pspace{\newpage}
    
\item 
  \begin{enumerate}
  \item
    \begin{problem}
    A bug is traveling along a straight line and its position at time $t$
    is given by $s(t)$. The differential equation 
    $\dfrac{ds}{dt}=3t^2-4t$ tells us about its \fitb{0.75in}{velocity.}
      When we see the differential equation $s' = 3t^2-4t$, we know we are
      looking for \cfitb{0.5in}{$s$} as a function of \cfitb{0.5in}{$t$}.
    \end{problem}

    \begin{solution}
        If $s(t)$ is the bug's position, then $\frac{ds}{dt}$ is the bug's
        velocity.
      The fact that the equation involves $s'$ tells us that $s$ is the
      function we're looking for; the fact that the other variable in the
      equation is $t$ tells us that $s$ must be a function of $t$. 
    \end{solution}

  \item

    \begin{problem}
      Which of the following are solutions of the differential equation
      $s' = 3t^2-4t$?
       \raggedcolumns
        \begin{enumerate} [label=\protect{\maybecircle{\Alph*.}}, itemsep=0.2cm]
        \plainitem $s(t)=t^3$
        \circleitem $s(t)=t^3-2t^2+3$
        \plainitem $s(t)=0$
        \end{enumerate}      
            
    \end{problem}

    \begin{solution}
      Let's look at each one:
      \begin{enumerate}[label=\Alph*.]
      \item If $s(t) = t^3$, then $s'(t) = 3t^2$,
      which is not the same as $3t^2
        - 4t$.
      \item If $s(t) = t^3 - 2t^2 + 3$, then $s'(t) = 3t^2 - 4t$.
      \item If $s(t) = 0$, then $s'(t) = 0$, which is not equal the same as
        $3t^2 - 4t$. 
      \end{enumerate}
    \end{solution}

  \item

    \note{Hopefully, students will see here that they can actually find
      the general solution by integrating.  Make sure they understand why
      this would not have worked in the previous problem.}

    \begin{problem}[\vfill]
      Can you find any more solutions of $s' = 3t^2 - 4t$?  Can you find
      them all?  (In other words, can you find the \uline{general
        solution} of $s' = 3t^2 - 4t$?)
    \end{problem}

    \begin{solution}
      The differential equation $s' = 3t^2 - 4t$ says that the derivative
      of $s$ is $3t^2 - 4t$, so the solutions are simply all
      antiderivatives of $3t^2 - 4t$; that is, $s(t) =
      \displaystyle \int (3t^2 - 4t)\,dt = \boxed{t^3 - 2t^2 + C}$. 
    \end{solution}

  \end{enumerate}
    

    
%   \item 
%     \begin{problem}[\vfill]
%       Let $P(t)$ be the population of turkeys in Cambridge. 
%       Can you think of any solutions of the differential equation
%       $\dfrac{dP}{dt} = 2t$?  Can you find all solutions?  (In other
%       words, can you find the general solution?)
%     \end{problem}

%     \begin{solution}
%       We're looking for functions $P(t)$ whose derivative is the function
%       $2t$; that is, we're looking for all antiderivatives of $2t$: $P(t)
%       = \displaystyle \int 2t\,dt = \boxed{t^2 + C}$.
%     \end{solution}


% \item 
%   The population of a bacteria colony increases at a rate proportional to
%   the population size.
%   \begin{enumerate}
%   \item
%     \begin{problem}[0.5in]
%       If $P(t)$ is the population size at time $t$, write a differential
%       equation that expresses the situation.      
%     \end{problem}

%     \begin{solution}
%       We can translate the information we're given into an equation: 
%       \[ \underbrace{\text{The rate at which the population
%           increases}}_{\frac{dP}{dt}} \text{ } \underbrace{\text{is
%           proportional to}}_{{} = k \text{ times}} \text{ }
%       \underbrace{\text{the population size}}_{P}.\] That is,
%       $\boxed{\frac{dP}{dt} = k P}$. 
%     \end{solution}

%   \item
%     \begin{problem}[\vfill]
%       When the colony has a population of 500 bacteria, the colony is
%       growing at a rate of 100 bacteria per hour.  Incorporate this into
%       your differential equation.
%     \end{problem}

%     \begin{solution}
%       We are told that, when $P = 500$, $\frac{dP}{dt} = 100$.  Plugging
%       this into our differential equation $\frac{dP}{dt} = kP$, $100 = 500
%       k$, so $k = 0.2$.  Therefore, our differential equation becomes
%       $\boxed{\frac{dP}{dt} = 0.2 P}$.
%     \end{solution}

  

%     \pspace{\newpage}

%   \item
%     \begin{problem}[\stretch{0.5}]
%       If the colony has an initial population of 200 bacteria, find a
%       formula for $P(t)$.
%     \end{problem}

%     \begin{solution}
%       The general solution to the differential equation $\frac{dP}{dt} =
%       0.2 P$ is $P(t) = C e^{0.2 t}$.  Knowing the initial population is
%       200 says that $P(0) = 200$, so $C e^{0.2 (0)} = 200$, which tells us
%       that $C = 200$.  So, $P(t) = \boxed{200 e^{0.2 t}}$.
%     \end{solution}


%   \end{enumerate}


\item 
  The mass of a radioactive substance decays at a rate proportional to the
  mass. 
  \begin{enumerate}
  \item
    \begin{problem}[0.25in]
      If $M(t)$ is the mass of the substance at time $t$, write a
      differential equation that expresses the situation.  How does the
      fact that the substance is decaying factor into your differential
      equation? 
    \end{problem}

    \begin{solution}
      We are told that
      \[ \underbrace{\text{The rate at which the mass
          decays}}_{\frac{dM}{dt}} \text{ } \underbrace{\text{is
          proportional to}}_{= k \text{ times}} \text{ }
      \underbrace{\text{the mass}}_{M}.\] That is, $\boxed{\frac{dM}{dt} =
        kM}$.  The fact that the substance is decaying tells us that
      $\frac{dM}{dt}$ is negative even though $M$ (the mass) is positive,
      so \fbox{the constant $k$ must be negative}.
    \end{solution}

  \item
    \begin{problem}[\vfill]
      If the mass is initially 10 g, and it is 6 g after 5 days, find a
      formula for $M(t)$.
    \end{problem}

    \begin{solution}
      Here are two possible approaches.
      \begin{itemize}
      \item \emph{Approach 1.}  The solutions of $\frac{dM}{dt} = kM$ are
        $M(t) = C e^{kt}$, and we can use the given information to find
        $C$ and $k$.
        \begin{align*}
          M(0) &= 10 \Longrightarrow C e^{k \cdot 0} = 10 \Longrightarrow
          C = 10 \\
          M(5) &= 6 \Longrightarrow C e^{k \cdot 5} = 6
        \end{align*}
        Since we've already figured out that $C = 10$, the second equation
        can be rewritten as
        \begin{align*}
          10 e^{5k} &= 6 \\
          e^{5k} &= 0.6 \\
          5k &= \ln (0.6) \\
          k &= \frac{\ln (0.6)}{5}
        \end{align*}
        So, $M(t) = \boxed{ 10 e^{\ln (0.6) t / 5}}$.

      \item \emph{Approach 2.}  The solutions of $\frac{dM}{dt} = kM$ are
        $M(t) = C e^{kt}$; that is, they are exponential functions.  So,
        we're looking for an exponential function $M(t)$ that satisfies
        $M(0) = 10$ and $M(5) = 6$.  Exponential functions have a constant
        percent change per unit time; this mass decays 40\% in the first 5
        days, so it must decay 40\% \emph{every} 5 days.  Therefore, a
        formula for $M(t)$ is $M(t) = \boxed{10 (0.6)^{t / 5}}$.

      \end{itemize}
    \end{solution}

  \item
    \begin{problem}[\stretch{0.5}]
      After 2.5 days, is the mass of the substance greater than, less
      than, or equal to 8 grams?  Explain how you can be sure. 
    \end{problem}

    \begin{solution}
      The substance decays from 10 grams to 6 grams in 5 days, so what
      happens in 2.5 days?  One way to approach this is to think
      graphically: the substance decays exponentially (shown in red
      below).  If it instead decayed linearly, then it would be 8 grams
      after 2.5 days (shown in blue below):
      \begin{center}
        \begin{tikzpicture}
          \begin{axis}[cycle list name=mycolor, cycle list shift=1,
            xtick={0, 2.5, 5}, ytick={6, 8, 10}]
            \addplot+[domain=0:10] {10 - 0.8*x};
            \addplot+[domain=0:10] {10 *(0.6)^(x/5)};
          \end{axis}
        \end{tikzpicture}
      \end{center}
      As we can see from the graph, the actual mass of the substance is
      \fbox{less than 8 grams} after 2.5 days. 
    \end{solution}

  \item
    \begin{problem}[\stretch{0.75}]
      What is the half-life of this substance?  (Remember that the
      half-life is the amount of time it takes for half of the substance
      to decay.)      
    \end{problem}

    \begin{solution}
      We'd like to figure out when the mass decays to be half of its
      initial size, i.e., when it's 5 grams.  In this solution, we'll use
      our second formula for $M(t)$, $M(t) = 10(0.6)^{t / 5}$ (you could
      use either, since they're equivalent):
      \begin{align*}
        10 (0.6)^{t / 5} &= 5 \\
        (0.6)^{t / 5} &= 0.5 \\
        \intertext{Taking the natural log of both sides,}
        \ln \left( (0.6)^{t / 5} \right) &= \ln (0.5) \\
        \frac{t}{5} \ln (0.6) &= \ln (0.5) \\
        5 &= \boxed{ \frac{5 \ln(0.5)}{\ln(0.6)}}
      \end{align*}
      This is about 6.8 days.
    \end{solution}

\pspace{\newpage}

  % \item

  %   \note{If you have time, this problem is an opportunity to start on
  %     some more advanced modeling.  (I will not do it unless I have a huge
  %     amount of time, since modeling can be confusing initially.  I have
  %     it here mostly for the students who tend to speed through the
  %     worksheet on their own.)}

  %   \begin{problem}[\stretch{0.5}]
  %     Suppose an environmental group is collecting this radioactive
  %     substance from a contaminated site in order to clean up the site.
  %     Each day, they collect 100 g of the substance and store it in a
  %     collection facility.  Let $C(t)$ be the mass of the substance in the
  %     collection facility $t$ days after they begin the cleanup.  Write a
  %     differential equation that models the situation. 
  %   \end{problem}

  %   \begin{solution}
  %     There are two factors causing the mass in the collection facility to
  %     change.  First, the group is collecting 100 g of the substance each
  %     day.  Second, the substance decays radioactively.

  %     When we write a differential equation, we are writing about rates of
  %     change.  Here, the rate of change of $C(t)$ is (rate at which the
  %     group collects the substance) $-$ (rate at which the substance
  %     decays).  The rate at which the group collects the substance is a
  %     constant rate of 100 g per day.  From
  %     {{{{\#5}}(a)}}, we can model the rate at
  %     which the substance decays as $k$ times the amount of the substance,
  %     or $k C$.  In {{{{\#5}}(b)}}, we found that
  %     the appropriate value of $k$ for this particular substance is $k =
  %     \frac{\ln (0.6)}{5}$.  So, putting this all together,
  %     \begin{align*}
  %       \frac{dC}{dt} &= \text{(rate at which the group collects the
  %         substance)} - \text{(rate at which the substance decays)} \\
  %       &= 100 - \frac{\ln (0.6)}{5} C
  %     \end{align*}
  %     So, our differential equation is $\boxed{\frac{dC}{dt} = 100 -
  %       \frac{\ln (0.6)}{5} C}$. 
  %   \end{solution}

  \end{enumerate}

% \item  % Andrew Dittmer

%   \begin{problem}
%     Which of the following are solutions to the differential equation
%     $\dfrac{d^2y}{dx^2} = 2x \left(\dfrac{dy}{dx} \right)^2$?
%   \end{problem}

  
%     \begin{enumerate}[label=\protect\maybecircle{\Alph*.}]
%     \circleitem $y = \frac{1}{x}$
%     \circleitem $y = 17$
%     \circleitem $y = -\arctan(x)$
%     \circleitem $y = \frac{1}{x} + 3$
%     \end{enumerate}
  

%   \pspace{\stretch{0.75}}

%   \begin{solution}
%     We can check each one:
%     \begin{enumerate}[label=\Alph*.]
%     \item If $y = \frac{1}{x} = x^{-1}$, then $\frac{dy}{dx} = -x^{-2}$,
%       so $\frac{d^2 y}{dx^2} = 2 x^{-3}$ and $2x \left( \frac{dy}{dx}
%       \right)^2 = 2x (-x^{-2})^2 = 2x (x^{-4}) = 2x^{-3}$.  So, $\frac{d^2
%         y}{dx^2} = 2x \left( \frac{dy}{dx} \right)^2$.
%     \item If $y = 17$, then $\frac{dy}{dx} = 0$, so both $\frac{d^2
%         y}{dx^2}$ and $2x \left( \frac{dy}{dx} \right)^2$ are equal to 0.
%     \item If $y = -\arctan (x)$, then $\frac{dy}{dx} = - \frac{1}{1 + x^2}
%       = - (1 + x^2)^{-1}$, so $\frac{d^2 y}{dx^2} = (1 + x^2)^{-2} (2x) =
%       \frac{2x}{(1 + x^2)^2}$ and $2x \left( \frac{dy}{dx} \right)^2 = 2x
%       \left( - \frac{1}{1 + x^2} \right)^2 = \frac{2x}{(1 + x^2)^2}$.  So,
%       $\frac{d^2 y}{dx^2} = 2x \left( \frac{dy}{dx} \right)^2$.
%     \item If $y = \frac{1}{x} + 3 = x^{-1} + 3$, then $\frac{dy}{dx} =
%       -x^{-2}$, so $\frac{d^2 y}{dx^2} = 2 x^{-3}$ and $2x \left(
%         \frac{dy}{dx} \right)^2 = 2x (-x^{-2})^2 = 2x (x^{-4}) = 2x^{-3}$.
%       So, $\frac{d^2 y}{dx^2} = 2x \left( \frac{dy}{dx} \right)^2$.
%     \end{enumerate}
%   \end{solution}

  \item 

  Let's look at the differential equation $\dfrac{dq}{dt} = -q^2$.

  \begin{enumerate}
  \item

    \begin{problem}[0.4in]
      This differential equation is telling us to look for
      \fitb{0.5in}{$q$} as a function of \fitb{0.5in}{$t$}.  Write the
      differential equation in words.
    \end{problem}

    \begin{solution}
      The fact that we're looking at a derivative $\frac{dq}{dt}$ tells us
      that we're thinking of $q$ as a function of $t$.

      In words, the equation $\dfrac{dq}{dt} = -q^2$ says, ``the
      derivative of the function $q(t)$ is equal to the negative of the
      function squared''.
    \end{solution}

  \item

    \begin{problem}
      Which of the following functions are solutions of the differential
      equation?
      
        \begin{enumerate}[label=\protect{\maybecircle{\Alph*.}}, itemsep=0.2cm]
        \plainitem $q(t) = - \frac{t^3}{3}$
        \circleitem $q(t) = \frac{1}{t}$
        \circleitem $q(t) = \frac{1}{t + 2}$
        \plainitem $q(t) = \frac{1}{2t + 1}$
        \plainitem $q(t) = \sqrt{t}$
        \circleitem $q(t) = 0$
        \end{enumerate}      
            
    \end{problem}

    \begin{solution}
      For each function $q(t)$, we'd like to determine whether the
      derivative of the function $q(t)$ is equal to the negative of the
      function squared (that's exactly what we said in
      {{\#1}}).  Let's look at each one:
      \begin{enumerate}[label=\Alph*.]
      \item If $q(t) = - \frac{t^3}{3}$, the derivative $\frac{dq}{dt}$ is
        equal to $-t^2$, while $-q^2$ is equal to $- \left( -
          \frac{t^3}{3} \right)^2 = - \frac{t^6}{9}$.  Therefore,
        $\frac{dq}{dt} \neq -q^2$.
      \item If $q(t) = \frac{1}{t}$, then $\frac{dq}{dt} = - t^{-2}$,
        while $-q^2 = - \left( \frac{1}{t} \right)^2 = - \frac{1}{t^2}$.
        So, in this case, $\frac{dq}{dt}$ is equal to $-q^2$.
      \item If $q(t) = \frac{1}{t + 2} = (t + 2)^{-1}$, then
        $\frac{dq}{dt} = - (t + 2)^{-2}$, while $-q^2$ is equal to $-
        \left( \frac{1}{t + 2} \right)^2$, so $\frac{dq}{dt}$ is equal to
        $-q^2$.
      \item If $q(t) = \frac{1}{2t + 1} = (2t + 1)^{-1}$, then
        $\frac{dq}{dt} = - (2t + 1)^{-2} \cdot 2 = - \frac{2}{(2t +
          1)^2}$, while $-q^2$ is equal to $-\left( \frac{1}{2t + 1}
        \right)^2 = - \frac{1}{(2t + 1)^2}$.  So, in this case,
        $\frac{dq}{dt}$ is not equal to $-q^2$.
      \item If $q(t) = \sqrt{t}$, then $\frac{dq}{dt} = \frac{1}{2}
        t^{-1/2}$, while $-q^2$ is equal to $- \left( \sqrt{t} \right)^2 =
        -t$.  So, in this case, $\frac{dq}{dt}$ is not equal to $-q^2$.
      \item If $q(t) = 0$, then $\frac{dq}{dt} = 0$, while $-q^2$ is equal
        to $0$, so $\frac{dq}{dt}$ is equal to $-q^2$. 
      \end{enumerate}
    \end{solution}

  \item

    \begin{problem}[0.75in]
      Can you guess any more solutions of the differential equation? 
    \end{problem}

    \begin{solution}
      Two of the solutions we saw were $q(t) = \frac{1}{t}$ and $q(t) =
      \frac{1}{t + 2}$, so we might guess that any function of the form
      $q(t) = \frac{1}{t \text{ + a constant}}$ is a solution of the
      differential equation. 

      It's easy to check whether our guess is correct, so let's do that.
      If we call the constant in our formula $C$, then $q(t) = \frac{1}{t
        + C}$.  The derivative $\frac{dq}{dt}$ is $-(t + C)^{-2} = -
      \frac{1}{(t + C)^2}$, while $-q^2 = - \left( \frac{1}{t + C}
      \right)^2 = - \frac{1}{(t + C)^2}$, so $\frac{dq}{dt}$ is equal to
      $-q^2$.
    \end{solution}

  \end{enumerate}

\end{enumerate}

\end{document}